% !TEX program = lualatex
\documentclass[12pt,aspectratio=43]{beamer}
% Shared Beamer design. Load after \documentclass using \input.
% Engine: LuaLaTeX. Fonts: UD Digi Kyokasho N (regular/bold), Arial.
\usepackage{luatexja-fontspec}
\setmainjfont[BoldFont={UD Digi Kyokasho N Bold}]{UD Digi Kyokasho N}
\setsansjfont[BoldFont={UD Digi Kyokasho N Bold}]{UD Digi Kyokasho N}
\setmonojfont[BoldFont={UD Digi Kyokasho N Bold}]{UD Digi Kyokasho N}
% Liberation Sans keeps the layout usable on systems without Arial.
\IfFontExistsTF{Arial}{
  \setmainfont{Arial}
  \setsansfont{Arial}
  \setmonofont{Arial}
}{
  \setmainfont{Liberation Sans}
  \setsansfont{Liberation Sans}
  \setmonofont{Liberation Sans}
}
\renewcommand{\kanjifamilydefault}{\gtdefault}

% Typeset mathematics with a dedicated TeX math font, independent of Arial.
\usefonttheme{professionalfonts}
\usepackage[mathrm=sym]{unicode-math}
\setmathfont{Latin Modern Math}
\setoperatorfont{\symup}

\definecolor{courseblue}{HTML}{4472C4}
\definecolor{ruleblue}{HTML}{8FAADC}
\setbeamercolor{normal text}{fg=black,bg=white}
\setbeamercolor{structure}{fg=courseblue}
\setbeamercolor{frametitle}{fg=black,bg=white}
\setbeamercolor{title banner}{fg=white,bg=courseblue}
\setbeamerfont{frametitle}{size=\fontsize{18}{22}\selectfont,series=\mdseries}
\setbeamerfont{quote}{shape=\upshape}
\setbeamersize{text margin left=9mm,text margin right=9mm}
\setbeamertemplate{navigation symbols}{}
\setbeamercolor{page number}{fg=gray,bg=white}
\setbeamertemplate{footline}{%
  \begin{beamercolorbox}[wd=\paperwidth,ht=2.5mm,dp=1.5mm]{page number}
    \hfill{\fontsize{8}{9}\selectfont
      \ifnum\value{framenumber}>0
        \insertframenumber\,/\,\inserttotalframenumber
      \fi}\hspace*{3mm}
  \end{beamercolorbox}}
\setbeamertemplate{headline}{%
  \hfill\textcolor{gray}{\fontsize{6}{7}\selectfont
    \hypersetup{linkcolor=gray}% Keep Beamer's title links in the header color.
    \ifnum\value{framenumber}=0 \insertshortdate\enspace\fi
    \insertshorttitle\space\insertshortsubtitle}\hspace{2mm}\vspace{1.5mm}}
\setbeamertemplate{frametitle}{%
  \vspace{2mm}\centering\insertframetitle\par
  \vspace{-1mm}{\color{ruleblue}\rule{.44\paperwidth}{.5pt}}\par
  \vspace{2mm}}
\setbeamercolor{itemize item}{fg=courseblue}
\setbeamercolor{itemize subitem}{fg=black}
\setbeamercolor{itemize subsubitem}{fg=black}
\setbeamertemplate{itemize item}{\textbullet}
\setbeamertemplate{itemize subitem}{\textnormal{-}}
\setbeamertemplate{itemize subsubitem}{\textnormal{-}}
\setlength{\parskip}{.45em}
\hypersetup{colorlinks=true,urlcolor=courseblue}

% Use \titlepage inside a [t,noframenumbering] frame. The background extends
% 1 mm past both page edges to avoid hairline gaps in PDF viewers.
\setbeamertemplate{title page}{%
  \vspace{2mm}
  \begin{beamercolorbox}[wd=\paperwidth,ht=43mm,dp=13mm,colsep*=1mm,center]{title banner}
    {\fontsize{32}{40}\selectfont\inserttitle\par}
    \vspace{2mm}
    {\fontsize{20}{24}\selectfont\insertsubtitle\par}
  \end{beamercolorbox}
  \vspace{10mm}
  \hfill\insertauthor
}

\setlength{\parskip}{0pt}
\usepackage{listings}
\usepackage{upquote}
\IfFontExistsTF{Consolas}{\newfontfamily\codefont{Consolas}}{\newfontfamily\codefont{Liberation Mono}}
\definecolor{codevalue}{HTML}{7A3E9D}
\definecolor{codebuiltin}{HTML}{006DB0}
\definecolor{codestring}{HTML}{267348}
\newcommand{\code}[1]{{\codefont #1}}

\lstdefinestyle{python}{
  language=Python,
  basicstyle=\codefont\fontsize{13}{15}\selectfont,
  keywordstyle=\color{codebuiltin},
  keywordstyle=[2]\color{codebuiltin},
  morekeywords=[2]{print},
  stringstyle=\color{codestring},
  showstringspaces=false,
  columns=fullflexible,
  keepspaces=true,
  upquote=true,
  tabsize=4,
  breaklines=true,
  breakatwhitespace=true,
  backgroundcolor=\color{black!4},
  frame=single,framerule=0pt,framesep=1.5mm,
  xleftmargin=1mm,xrightmargin=1mm,
  aboveskip=1mm,belowskip=0.5mm,
  literate=*{0}{{{\color{codevalue}0}}}1
    {1}{{{\color{codevalue}1}}}1 {2}{{{\color{codevalue}2}}}1
    {3}{{{\color{codevalue}3}}}1 {4}{{{\color{codevalue}4}}}1
    {5}{{{\color{codevalue}5}}}1 {6}{{{\color{codevalue}6}}}1
    {7}{{{\color{codevalue}7}}}1 {8}{{{\color{codevalue}8}}}1
    {9}{{{\color{codevalue}9}}}1
}
\lstset{style=python}
\newcommand{\out}[1]{%
  \par\vspace{.5mm}\noindent
  {\small\textcolor{black!60}{出力}}\hspace{1mm}%
  \begin{minipage}[t]{.88\linewidth}\vspace{0pt}%
    {\codefont\fontsize{11}{12.5}\selectfont #1}%
  \end{minipage}\par}
\setbeamercolor{lesson note}{fg=black,bg=courseblue!9}
\newenvironment{lessonbox}
  {\begin{beamercolorbox}[wd=\linewidth,sep=1.5mm]{lesson note}}
  {\end{beamercolorbox}\par}

\title[計算物理学B]{計算物理学B}
\subtitle[第四回付録]{第四回付録：arrayとndarray}
\author{野垣 康介、藤本 悠輝}
\date[2026/11/10]{2026/11/10}
\hypersetup{pdftitle={計算物理学B 第四回付録：arrayとndarray},pdfauthor={野垣 康介、藤本 悠輝}}

\begin{document}

% 表紙
\begin{frame}[t,noframenumbering]
  \titlepage
\end{frame}

% 1. 名前と基本的な違い
\begin{frame}[t]{3種類の配列：名前と役割}
  \small
  \vspace{-1mm}
  \begingroup
    \renewcommand{\arraystretch}{1.12}
    \begin{tabular}{@{}p{3.3cm}p{6.6cm}@{}}
      \hline
      \textbf{名前} & \textbf{できること・特徴} \\
      \hline
      \code{list} &
        Python組込み。異なる型を混ぜられ、要素を追加・削除できます。 \\
      \code{array.array} &
        標準ライブラリの可変1次元配列。型コードで要素型をそろえます。 \\
      \mbox{\code{numpy.ndarray}} &
        NumPyの1次元・多次元配列。\code{dtype}と\code{shape}を持ち、要素ごとの計算や行列計算に便利です。 \\
      \hline
    \end{tabular}\par
  \endgroup
  \vspace{1mm}
  \begin{lessonbox}
    \code{np.array(...)}は作成関数、\code{ndarray}は配列の型です。1次元の結果も\code{ndarray}です。
  \end{lessonbox}
\end{frame}

% 2. 作り方と型の指定
\begin{frame}[t,fragile]{作り方と型の指定}
  \vspace{-1mm}
\begin{lstlisting}
import array as ar
import numpy as np
a_std = ar.array("d", [1.0, 2.0])
a_np = np.array([1.0, 2.0], dtype=float)
print(a_std.typecode, a_np.dtype, a_np.shape)
\end{lstlisting}
  \out{d float64 (2,)}
  \small \code{'d'}はCのdouble、\code{'i'}はCのint（幅は環境依存）で、型コードはPythonの版にも依存します。\code{np.array(...)}は型を推定し、\code{dtype}で明示もできます。\code{typecode}とは別の指定方式です。
\end{frame}

% 3. 演算子の違い
\begin{frame}[t,fragile]{同じ演算子でも意味が変わる}
  \vspace{-1mm}
\begin{lstlisting}
a_std = ar.array("d", [1.0, 2.0])
a_np = np.array([1.0, 2.0])
print((a_std * 2).tolist())
print(a_np * 2)
\end{lstlisting}
  \out{\begin{tabular}{@{}l@{}}
    [1.0, 2.0, 1.0, 2.0]\\\relax
    [2. 4.]
  \end{tabular}}
  \vspace{-1mm}
  \begin{lessonbox}
    同じ型コードの標準\code{array}の\code{+}は連結、\code{* 2}は繰り返し。\code{**}は\code{TypeError}です。\\
    \code{ndarray}の\code{+}、\code{*}、\code{**}は要素ごとの演算、\code{@}は行列積です。
  \end{lessonbox}
\end{frame}

% 4. 標準arrayの追加・削除
\begin{frame}[t,fragile]{標準\code{array}の追加・削除}
  \begin{columns}[T,onlytextwidth]
    \begin{column}{.48\textwidth}
      \small すべて元配列を書き換えます。
      \par\vspace{1mm}
      \code{append(x)}・\code{extend(xs)}：末尾へ追加。\\
      \code{insert(i, x)}：添字位置へ追加。\\
      \code{del a[i]}：その要素を削除。\\
      \code{pop(i)}：削除して値を返す。\code{i}省略時は末尾。\\
      \code{remove(x)}：最初に一致する値を1つ削除。
    \end{column}
    \begin{column}{.51\textwidth}
\begin{lstlisting}
a = ar.array("d", [1, 2, 2, 3])
a.append(4)
print(a.pop())
a.remove(2)
print(a.tolist())
\end{lstlisting}
      \out{4.0\par [1.0, 2.0, 3.0]}
    \end{column}
  \end{columns}
  \vspace{1mm}
  \small 範囲外・空の\code{pop}は\code{IndexError}、見つからない値への\code{remove}は\code{ValueError}です。
\end{frame}

% 5. NumPyの追加
\begin{frame}[t,fragile]{NumPyの追加は返り値を受け取る}
  \small \code{ndarray}に\code{append}メソッドはありません。\code{np.append}・\code{np.insert}は新配列を返します。
\begin{lstlisting}
a_np = np.array([1.0, 2.0, 2.0, 3.0])
grown_np = np.append(a_np, 4)
inserted_np = np.insert(a_np, 1, 9)
print(a_np, grown_np, inserted_np, sep="\n")
\end{lstlisting}
  \out{\begin{tabular}{@{}l@{}}
    [1. 2. 2. 3.]\\\relax
    [1. 2. 2. 3. 4.]\\\relax
    [1. 9. 2. 2. 3.]
  \end{tabular}}
  \small 元に反映するには返り値を代入します。
\end{frame}

% 6. NumPyの削除
\begin{frame}[t,fragile]{NumPyの削除：値か、添字か}
  \small \code{np.delete}は添字を指定し、条件式は条件に合う値を選びます。
\begin{lstlisting}
a_np = np.array([1.0, 2.0, 2.0, 3.0])
print(np.delete(a_np, 1))
print(a_np[a_np != 2])
print(a_np)
\end{lstlisting}
  \out{\begin{tabular}{@{}l@{}}
    [1. 2. 3.]\\\relax
    [1. 3.]\\\relax
    [1. 2. 2. 3.]
  \end{tabular}}
  \small \code{np.delete(a\_np, 1)}は添字1を削除。\code{a\_np[a\_np != 2]}は値2をすべて除きます。どちらも新配列を返します。直接の添字代入は元を更新します。
\end{frame}

% 7. axis
\begin{frame}[t,fragile]{多次元配列の追加と\code{axis}}
\begin{lstlisting}
m = np.array([[1, 0], [0, 2]])
r = np.array([[-1, 0]])
joined = np.concatenate((m, r), axis=0)
flat = np.append(m, r)
print(joined.shape, flat.shape)
\end{lstlisting}
  \out{(3, 2) (6,)}
  \small \code{axis=0}は行、\code{axis=1}は列。結合する軸以外の大きさをそろえます。\\
  \code{np.append}・\code{np.insert}・\code{np.delete}は既定の\code{axis=None}で平坦化。\code{np.concatenate}は既定の\code{axis=0}です。
\end{frame}

% 8. スライスとコピー
\begin{frame}[t,fragile]{スライスはコピーか、共有か}
\begin{lstlisting}
a_std = ar.array("d", [1, 2, 3, 4])
s_std = a_std[1:3]
s_std[0] = 99
a_np = np.array([1., 2., 3., 4.])
s_np = a_np[1:3]
s_np[0] = 99
print(a_std.tolist(), a_np, sep="\n")
\end{lstlisting}
  \out{\begin{tabular}{@{}l@{}}
    [1.0, 2.0, 3.0, 4.0]\\\relax
    [ 1. 99.  3.  4.]
  \end{tabular}}
\end{frame}

% 9. viewとcopy・高度な添字
\begin{frame}[t,fragile]{viewとcopy・高度な添字}
\begin{lstlisting}
a_np = np.array([1., 2., 3., 4.])
selected_np = a_np[[1, 2]]
selected_np[0] = 99
copy_np = a_np[1:3].copy()
copy_np[0] = -1
print(a_np)
\end{lstlisting}
  \out{[1. 2. 3. 4.]}
  \small 標準\code{array}のスライスはコピー、NumPyの基本スライスはviewです。整数配列・条件配列による選択結果と\code{.copy()}も独立配列です。直接の添字代入は元を更新し、\code{b = a}は同じオブジェクトを参照します。
\end{frame}

% 10. shape, reshape, 転置と行列積
\begin{frame}[t,fragile]{\code{shape}・\code{reshape}・転置・行列積}
\begin{lstlisting}
matrix_np = np.arange(4).reshape(2, 2)
vector_np = np.array([1, 2])
print((matrix_np @ vector_np).tolist())
print(matrix_np.T.tolist())
print(vector_np.shape, vector_np.T.shape)
\end{lstlisting}
  \out{\begin{tabular}{@{}l@{}}
    [2, 8]\\\relax
    [[0, 2], [1, 3]]\\
    (2,) (2,)
  \end{tabular}}
  \small \code{reshape}は要素数を保ち、結果はviewまたはcopyです。\code{@}は行列積。\code{.T}は2次元では転置、1次元では形を変えません。
\end{frame}

% 11. 標準arrayでの整数代入
\begin{frame}[t,fragile]{整数型への小数代入：標準\code{array}}
  \small \code{array('i')}の要素はCのint型です。小数はそのまま代入できません。
\begin{lstlisting}
int_std = ar.array("i", [1, 2])
int_std[0] = -2.5
\end{lstlisting}
  \out{\code{TypeError}}
\end{frame}

% 12. NumPyでの整数代入
\begin{frame}[t,fragile]{整数\code{ndarray}へ小数を代入すると}
\begin{lstlisting}
int_np = np.array([1, 2], dtype=int)
int_np[0] = -2.5
float_np = np.array([1, 2],
                    dtype=float)
float_np[0] = -2.5
print(int_np, float_np)
\end{lstlisting}
  \out{[-2  2] [-2.5  2. ]}
  \small この代入では小数部が0方向に落ち、\code{-2}になります。小数を保持するなら\code{dtype=float}で作ります。これはこの要素代入例の結果です。
\end{frame}

% 13. 相互変換とコピー
\begin{frame}[t,fragile]{相互変換とコピー}
\vspace{-1mm}
\begin{lstlisting}
source_std = ar.array("d", [0, 1])
converted_np = np.array(source_std)
converted_np[0] = 99
print(source_std.tolist(), converted_np)
print(np.sin(source_std), type(np.sin(source_std)))
\end{lstlisting}
  \out{\begin{tabular}{@{}l@{}}
    [0.0, 1.0] [99.  1.]\\\relax
    [0.         0.84147098] \textless class 'numpy.ndarray'\textgreater
  \end{tabular}}
  \small \code{np.sin}は標準\code{array}を受け取り\code{ndarray}を返します。共有・コピー条件は変換方法で異なります。
\end{frame}

% 14. 2次元配列を標準arrayへ戻す
\begin{frame}[t,fragile]{2次元配列を標準\code{array}へ戻す}
  \small 1次元なら\code{ar.array('d', nd.tolist())}で戻せます。2次元の\code{tolist()}は入れ子なので、平坦化します。
\begin{lstlisting}
matrix_np = np.array([[1.0, 2.0], [3.0, 4.0]])
flat_std = ar.array("d", matrix_np.flatten().tolist())
print(flat_std.tolist())
\end{lstlisting}
  \out{[1.0, 2.0, 3.0, 4.0]}
  \small \code{flatten()}後は値が並ぶだけで、\code{(2, 2)}の形は残りません。形が必要なら\code{ndarray}で扱います。
\end{frame}

% 15. 物理の計算：標準array
\begin{frame}[t,fragile]{自由落下を標準\code{array}で計算する}
  \small 静止・原点から落下、下向きを正、空気抵抗なし、\(g=10.0\ \mathrm{m/s^2}\)とします。
  \(\displaystyle x(t)=\frac{1}{2}gt^2=5.0t^2\quad(t\geq0)\)
\begin{lstlisting}
times_std = ar.array("d", [0., 1., 2.])
positions_std = ar.array("d")
for time_value in times_std:
    positions_std.append(
        5.0 * time_value ** 2)
print(positions_std.tolist())
\end{lstlisting}
  \out{[0.0, 5.0, 20.0]}
  \small 時刻を順番に走査し、各時刻の位置を\code{append}で追加します。
\end{frame}

% 16. 物理の計算：ndarray
\begin{frame}[t,fragile]{同じ自由落下を\code{ndarray}で計算する}
  \small 同じ条件・時刻では、配列演算で全位置をまとめて求められます。
\begin{lstlisting}
times_np = np.array([0., 1., 2.])
positions_np = 5.0 * times_np ** 2
print(positions_np)
\end{lstlisting}
  \out{[ 0.  5. 20.]}
  \begin{lessonbox}
    1次元の同型値を標準ライブラリだけで扱うなら\code{array.array}、一括計算・多次元・行列計算なら\code{ndarray}が便利です。速度は実際の条件で測ります。
  \end{lessonbox}
  \vfill
  {\fontsize{8}{10}\selectfont
    仕様・使用例：
    \href{https://docs.python.org/3/library/array.html}{Python：array}、
    \href{https://numpy.org/doc/stable/user/quickstart.html}{NumPy：基本操作}、
    \href{https://numpy.org/doc/stable/user/basics.copies.html}{コピーとview}\\
    \href{https://numpy.org/doc/stable/reference/generated/numpy.append.html}{np.append}、
    \href{https://numpy.org/doc/stable/reference/generated/numpy.delete.html}{np.delete}、
    \href{https://numpy.org/doc/stable/reference/generated/numpy.array.html}{np.array}}
\end{frame}

\end{document}
